% Options for packages loaded elsewhere
\PassOptionsToPackage{unicode}{hyperref}
\PassOptionsToPackage{hyphens}{url}
%
\documentclass[
]{book}
\usepackage{lmodern}
\usepackage{amssymb,amsmath}
\usepackage{ifxetex,ifluatex}
\ifnum 0\ifxetex 1\fi\ifluatex 1\fi=0 % if pdftex
  \usepackage[T1]{fontenc}
  \usepackage[utf8]{inputenc}
  \usepackage{textcomp} % provide euro and other symbols
\else % if luatex or xetex
  \usepackage{unicode-math}
  \defaultfontfeatures{Scale=MatchLowercase}
  \defaultfontfeatures[\rmfamily]{Ligatures=TeX,Scale=1}
\fi
% Use upquote if available, for straight quotes in verbatim environments
\IfFileExists{upquote.sty}{\usepackage{upquote}}{}
\IfFileExists{microtype.sty}{% use microtype if available
  \usepackage[]{microtype}
  \UseMicrotypeSet[protrusion]{basicmath} % disable protrusion for tt fonts
}{}
\makeatletter
\@ifundefined{KOMAClassName}{% if non-KOMA class
  \IfFileExists{parskip.sty}{%
    \usepackage{parskip}
  }{% else
    \setlength{\parindent}{0pt}
    \setlength{\parskip}{6pt plus 2pt minus 1pt}}
}{% if KOMA class
  \KOMAoptions{parskip=half}}
\makeatother
\usepackage{xcolor}
\IfFileExists{xurl.sty}{\usepackage{xurl}}{} % add URL line breaks if available
\IfFileExists{bookmark.sty}{\usepackage{bookmark}}{\usepackage{hyperref}}
\hypersetup{
  pdftitle={Data Science Methods for Cognitive Science},
  pdfauthor={Dan Hicks, hicks.daniel.j@gmail.com},
  hidelinks,
  pdfcreator={LaTeX via pandoc}}
\urlstyle{same} % disable monospaced font for URLs
\usepackage{color}
\usepackage{fancyvrb}
\newcommand{\VerbBar}{|}
\newcommand{\VERB}{\Verb[commandchars=\\\{\}]}
\DefineVerbatimEnvironment{Highlighting}{Verbatim}{commandchars=\\\{\}}
% Add ',fontsize=\small' for more characters per line
\usepackage{framed}
\definecolor{shadecolor}{RGB}{248,248,248}
\newenvironment{Shaded}{\begin{snugshade}}{\end{snugshade}}
\newcommand{\AlertTok}[1]{\textcolor[rgb]{0.94,0.16,0.16}{#1}}
\newcommand{\AnnotationTok}[1]{\textcolor[rgb]{0.56,0.35,0.01}{\textbf{\textit{#1}}}}
\newcommand{\AttributeTok}[1]{\textcolor[rgb]{0.77,0.63,0.00}{#1}}
\newcommand{\BaseNTok}[1]{\textcolor[rgb]{0.00,0.00,0.81}{#1}}
\newcommand{\BuiltInTok}[1]{#1}
\newcommand{\CharTok}[1]{\textcolor[rgb]{0.31,0.60,0.02}{#1}}
\newcommand{\CommentTok}[1]{\textcolor[rgb]{0.56,0.35,0.01}{\textit{#1}}}
\newcommand{\CommentVarTok}[1]{\textcolor[rgb]{0.56,0.35,0.01}{\textbf{\textit{#1}}}}
\newcommand{\ConstantTok}[1]{\textcolor[rgb]{0.00,0.00,0.00}{#1}}
\newcommand{\ControlFlowTok}[1]{\textcolor[rgb]{0.13,0.29,0.53}{\textbf{#1}}}
\newcommand{\DataTypeTok}[1]{\textcolor[rgb]{0.13,0.29,0.53}{#1}}
\newcommand{\DecValTok}[1]{\textcolor[rgb]{0.00,0.00,0.81}{#1}}
\newcommand{\DocumentationTok}[1]{\textcolor[rgb]{0.56,0.35,0.01}{\textbf{\textit{#1}}}}
\newcommand{\ErrorTok}[1]{\textcolor[rgb]{0.64,0.00,0.00}{\textbf{#1}}}
\newcommand{\ExtensionTok}[1]{#1}
\newcommand{\FloatTok}[1]{\textcolor[rgb]{0.00,0.00,0.81}{#1}}
\newcommand{\FunctionTok}[1]{\textcolor[rgb]{0.00,0.00,0.00}{#1}}
\newcommand{\ImportTok}[1]{#1}
\newcommand{\InformationTok}[1]{\textcolor[rgb]{0.56,0.35,0.01}{\textbf{\textit{#1}}}}
\newcommand{\KeywordTok}[1]{\textcolor[rgb]{0.13,0.29,0.53}{\textbf{#1}}}
\newcommand{\NormalTok}[1]{#1}
\newcommand{\OperatorTok}[1]{\textcolor[rgb]{0.81,0.36,0.00}{\textbf{#1}}}
\newcommand{\OtherTok}[1]{\textcolor[rgb]{0.56,0.35,0.01}{#1}}
\newcommand{\PreprocessorTok}[1]{\textcolor[rgb]{0.56,0.35,0.01}{\textit{#1}}}
\newcommand{\RegionMarkerTok}[1]{#1}
\newcommand{\SpecialCharTok}[1]{\textcolor[rgb]{0.00,0.00,0.00}{#1}}
\newcommand{\SpecialStringTok}[1]{\textcolor[rgb]{0.31,0.60,0.02}{#1}}
\newcommand{\StringTok}[1]{\textcolor[rgb]{0.31,0.60,0.02}{#1}}
\newcommand{\VariableTok}[1]{\textcolor[rgb]{0.00,0.00,0.00}{#1}}
\newcommand{\VerbatimStringTok}[1]{\textcolor[rgb]{0.31,0.60,0.02}{#1}}
\newcommand{\WarningTok}[1]{\textcolor[rgb]{0.56,0.35,0.01}{\textbf{\textit{#1}}}}
\usepackage{longtable,booktabs}
% Correct order of tables after \paragraph or \subparagraph
\usepackage{etoolbox}
\makeatletter
\patchcmd\longtable{\par}{\if@noskipsec\mbox{}\fi\par}{}{}
\makeatother
% Allow footnotes in longtable head/foot
\IfFileExists{footnotehyper.sty}{\usepackage{footnotehyper}}{\usepackage{footnote}}
\makesavenoteenv{longtable}
\usepackage{graphicx}
\makeatletter
\def\maxwidth{\ifdim\Gin@nat@width>\linewidth\linewidth\else\Gin@nat@width\fi}
\def\maxheight{\ifdim\Gin@nat@height>\textheight\textheight\else\Gin@nat@height\fi}
\makeatother
% Scale images if necessary, so that they will not overflow the page
% margins by default, and it is still possible to overwrite the defaults
% using explicit options in \includegraphics[width, height, ...]{}
\setkeys{Gin}{width=\maxwidth,height=\maxheight,keepaspectratio}
% Set default figure placement to htbp
\makeatletter
\def\fps@figure{htbp}
\makeatother
\setlength{\emergencystretch}{3em} % prevent overfull lines
\providecommand{\tightlist}{%
  \setlength{\itemsep}{0pt}\setlength{\parskip}{0pt}}
\setcounter{secnumdepth}{5}
\usepackage{booktabs}
\usepackage[]{biblatex}
\addbibresource{book.bib}
\addbibresource{packages.bib}

\title{Data Science Methods for Cognitive Science}
\author{Dan Hicks, \href{mailto:hicks.daniel.j@gmail.com}{\nolinkurl{hicks.daniel.j@gmail.com}}}
\date{2020-07-09}

\begin{document}
\maketitle

{
\setcounter{tocdepth}{1}
\tableofcontents
}
\hypertarget{preface}{%
\chapter*{Preface}\label{preface}}
\addcontentsline{toc}{chapter}{Preface}

This document is a set of notes (and perhaps later a textbook) for a course on data science methods for graduate students in social and behavioral sciences, taught by Professor Dan Hicks (they/them), UC Merced.

\hypertarget{week-0-setting-up}{%
\chapter*{Week 0: Setting up}\label{week-0-setting-up}}
\addcontentsline{toc}{chapter}{Week 0: Setting up}

Before the first day of class, you'll need to set up three things: \textbf{R, RStudio, and a GitHub account}. R is the primary piece of software we'll be learning to use. RStudio is an IDE (integrated development environment) for using R. (Think of R as the engine of a car and RStudio as the dashboard, steering wheel, and pedals.) GitHub is a service for sharing and collaborating on software engineering.

I'll assume you have R and RStudio installed using either \emph{the cloud way} or \emph{the Mac way}. If your computer is running Windows, Linux, or some other operating system that isn't Mac OS, I'm going to assume you have things set up the cloud way. These other operating systems are perfectly fine. But I don't use them, which makes it much harder for me to help you if things aren't working. And Windows handles things like file paths and text file formats in idiosyncratic ways that can cause weird errors.

In either case, all of the software and services we'll be using are either free or freemium. You won't need any of the paid extras for this course.

\hypertarget{r-and-rstudio-the-mac-way}{%
\section{R and RStudio: The Mac Way}\label{r-and-rstudio-the-mac-way}}

With this way, R and RStudio are installed and run on your computer, and all of the data files that you work with will also be on your computer. Unless you're running a Chromebook or similar, your computer is almost certainly more powerful than the default instance that RStudio Cloud gives you (see below), and you can use familiar tools to do things like organize folders and preview files. But sometimes particular details about the way your computer is set up can cause R and RStudio to do weird things, and these problems can be very difficult to diagnose and solve.

\begin{enumerate}
\def\labelenumi{\arabic{enumi}.}
\item
  Go to \url{https://cran.r-project.org/}. Download and install the latest release of R for your operating system (which should be Mac OS if you're following this way). As of 2020-06-02, the latest release is 4.0.0.
\item
  Go to \href{https://rstudio.com/products/rstudio/download/\#download}{\texttt{rstudio.com}}. Download and install the latest release of RStudio Desktop for your operating system (which should be Mac OS if you're following this way). As of 2020-06-02, the latest release is 1.3.959.
\end{enumerate}

\hypertarget{r-and-rstudio-the-cloud-way}{%
\section{R and RStudio: The Cloud Way}\label{r-and-rstudio-the-cloud-way}}

(\textbf{Note 2017-07-07} Starting in August, RStudio Cloud is moving out of beta and becoming a regular subscription service. Information on pricing is at \url{https://rstudio.cloud/plans/free}. We'll have to try it out, but it seems likely that the Free tier won't give you enough compute time (``project hours'') for this course. I can subscribe to the Instructor tier, but that has its own limitations. We'll have to see how things go.)

With this way, R and RStudio are installed and run on virtual computers (called \emph{instances}) on a big server somewhere. Each project can have its own instance, meaning that it runs on a nice clean computer and shouldn't have any weirdness. The RStudio Cloud service is currently in beta; it's free to use, but not completely guaranteed to be free of bugs. In order to use the service, you have to be connected to the internet, and you'll have to use the RStudio interface or a terminal window to do things like move files around. Also, the instances aren't very powerful, and so they'll be slow if you try to work with large data sets or do a complex analysis.

\begin{enumerate}
\def\labelenumi{\arabic{enumi}.}
\tightlist
\item
  Go to \href{https://rstudio.cloud/}{\texttt{rstudio.cloud}}. Look for ``Sign Up'' in the upper right.\footnote{RStudio gives you the option to log in with Google or GitHub. As a general practice, I recommend always keeping your accounts separate. Use a password manager to generate separate, strong passwords for each site. I use \href{https://www.lastpass.com/}{LastPass}.}
\end{enumerate}

\hypertarget{github}{%
\section{GitHub}\label{github}}

\begin{enumerate}
\def\labelenumi{\arabic{enumi}.}
\item
  Go to \href{https://github.com/}{\texttt{github.com}}. You can create a new account right on the front page.
\item
  You can fill in your profile if you want.
\end{enumerate}

\hypertarget{screen-layouts}{%
\section{Screen Layouts}\label{screen-layouts}}

In this course, lectures will be based around \emph{live coding}. This means that I (or whoever's doing the lecturing) will be writing code in real time, and you'll follow along, transcribing the code into RStudio (or whatever) on your local machine and running it at the same time as me.

Live coding is a great way to learn coding practices. But it requires you to have quick access to 3-5 different things:

\begin{itemize}
\tightlist
\item
  My screenshare (so you can see what I'm typing)
\item
  The course chat (so you can ask for help or other questions, or make comments)
\item
  Possibly a separate way to put up a flag so that I can see immediately that you have a question
\item
  Your own workspace (RStudio or whatever)
\item
  Possibly a separate app for taking notes
\end{itemize}

In an ordinary classroom, the first three things would all be handled by the classroom environment. In a fully online setting, everything needs to be handled on your computer.

The most effective way to do this is to have multiple displays. This is also the most expensive.

\href{https://carpentries.org/blog/2020/06/online-workshop-logistics-and_screen-layouts/}{This blog post} discusses alternatives to multiple displays. In particular, \href{https://carpentries.org/images//blog/2020/06/screen_layouts/02-student-with_participantlist-1366x768.jpg}{this window layout} fits the screenshare, chat, and your workspace on to a single laptop display. You can also have a separate text file open in RStudio for notes.

\hypertarget{introduction}{%
\chapter{Introduction}\label{introduction}}

\hypertarget{what-this-course-isnt-and-is}{%
\section{What this course isn't, and is}\label{what-this-course-isnt-and-is}}

Is Not:

\begin{itemize}
\tightlist
\item
  statistics (in the way you think)
\item
  general introduction to software engineering
\end{itemize}

Is:

\begin{itemize}
\tightlist
\item
  data cleaning, data management, and exploratory data analysis
\item
  habituation to some good software engineering practices that are especially valuable for data science work
\end{itemize}

\hypertarget{notes-on-some-design-decisions}{%
\section{Notes on some design decisions}\label{notes-on-some-design-decisions}}

\textbf{R, not Python}

\begin{itemize}
\tightlist
\item
  Python is a perfectly good language for general programming

  \begin{itemize}
  \tightlist
  \item
    And has some advantages over R in certain data science aspects, including working with strings and developing custom data structures
  \end{itemize}
\item
  \href{https://blog.shotwell.ca/posts/why_i_use_r/}{R is a better language for doing data science}

  \begin{itemize}
  \tightlist
  \item
    Vectors, data frames, and common statistical models are installed and loaded by default
  \item
    R is more \href{https://towardsdatascience.com/python-vs-r-choosing-the-best-programming-languages-for-data-science-b1327f01f6bf}{functional}
  \end{itemize}
\item
  R is also better for beginners

  \begin{itemize}
  \tightlist
  \item
    CRAN
  \item
    RStudio
  \end{itemize}
\end{itemize}

\textbf{Scripts, not notebooks}

\begin{itemize}
\item
  Quite a few scientists like the way \href{https://www.nature.com/articles/d41586-018-07196-1}{notebooks mix together code, visualizations, tables, and text}
\item
  \href{https://towardsdatascience.com/version-control-with-jupyter-notebooks-f096f4d7035a}{Notebooks do not play nicely with version control}
\item
  \href{https://docs.google.com/presentation/d/1n2RlMdmv1p25Xy5thJUhkKGvjtV-dkAIsUXP-AL4ffI/edit?usp=sharing}{Notebooks encourage bad habits and discourage good ones}
\end{itemize}

\hypertarget{lets-get-started-with-rstudio}{%
\section{Let's get started with RStudio}\label{lets-get-started-with-rstudio}}

\begin{itemize}
\tightlist
\item
  Quick tour of RStudio's UI
\item
  Let's plot a thing
\end{itemize}

\begin{Shaded}
\begin{Highlighting}[]
\KeywordTok{data}\NormalTok{(anscombe) }\CommentTok{\#\# \textless{}https://en.wikipedia.org/wiki/Anscombe\%27s\_quartet\textgreater{}}

\NormalTok{anscombe}
\end{Highlighting}
\end{Shaded}

\begin{verbatim}
##    x1 x2 x3 x4    y1   y2    y3    y4
## 1  10 10 10  8  8.04 9.14  7.46  6.58
## 2   8  8  8  8  6.95 8.14  6.77  5.76
## 3  13 13 13  8  7.58 8.74 12.74  7.71
## 4   9  9  9  8  8.81 8.77  7.11  8.84
## 5  11 11 11  8  8.33 9.26  7.81  8.47
## 6  14 14 14  8  9.96 8.10  8.84  7.04
## 7   6  6  6  8  7.24 6.13  6.08  5.25
## 8   4  4  4 19  4.26 3.10  5.39 12.50
## 9  12 12 12  8 10.84 9.13  8.15  5.56
## 10  7  7  7  8  4.82 7.26  6.42  7.91
## 11  5  5  5  8  5.68 4.74  5.73  6.89
\end{verbatim}

\begin{Shaded}
\begin{Highlighting}[]
\KeywordTok{str}\NormalTok{(anscombe)}
\end{Highlighting}
\end{Shaded}

\begin{verbatim}
## 'data.frame':	11 obs. of  8 variables:
##  $ x1: num  10 8 13 9 11 14 6 4 12 7 ...
##  $ x2: num  10 8 13 9 11 14 6 4 12 7 ...
##  $ x3: num  10 8 13 9 11 14 6 4 12 7 ...
##  $ x4: num  8 8 8 8 8 8 8 19 8 8 ...
##  $ y1: num  8.04 6.95 7.58 8.81 8.33 ...
##  $ y2: num  9.14 8.14 8.74 8.77 9.26 8.1 6.13 3.1 9.13 7.26 ...
##  $ y3: num  7.46 6.77 12.74 7.11 7.81 ...
##  $ y4: num  6.58 5.76 7.71 8.84 8.47 7.04 5.25 12.5 5.56 7.91 ...
\end{verbatim}

\begin{Shaded}
\begin{Highlighting}[]
\KeywordTok{summary}\NormalTok{(anscombe)}
\end{Highlighting}
\end{Shaded}

\begin{verbatim}
##        x1             x2             x3             x4           y1        
##  Min.   : 4.0   Min.   : 4.0   Min.   : 4.0   Min.   : 8   Min.   : 4.260  
##  1st Qu.: 6.5   1st Qu.: 6.5   1st Qu.: 6.5   1st Qu.: 8   1st Qu.: 6.315  
##  Median : 9.0   Median : 9.0   Median : 9.0   Median : 8   Median : 7.580  
##  Mean   : 9.0   Mean   : 9.0   Mean   : 9.0   Mean   : 9   Mean   : 7.501  
##  3rd Qu.:11.5   3rd Qu.:11.5   3rd Qu.:11.5   3rd Qu.: 8   3rd Qu.: 8.570  
##  Max.   :14.0   Max.   :14.0   Max.   :14.0   Max.   :19   Max.   :10.840  
##        y2              y3              y4        
##  Min.   :3.100   Min.   : 5.39   Min.   : 5.250  
##  1st Qu.:6.695   1st Qu.: 6.25   1st Qu.: 6.170  
##  Median :8.140   Median : 7.11   Median : 7.040  
##  Mean   :7.501   Mean   : 7.50   Mean   : 7.501  
##  3rd Qu.:8.950   3rd Qu.: 7.98   3rd Qu.: 8.190  
##  Max.   :9.260   Max.   :12.74   Max.   :12.500
\end{verbatim}

\begin{Shaded}
\begin{Highlighting}[]
\KeywordTok{library}\NormalTok{(ggplot2)}
\KeywordTok{ggplot}\NormalTok{(}\DataTypeTok{data =}\NormalTok{ anscombe, }\KeywordTok{aes}\NormalTok{(}\DataTypeTok{x =}\NormalTok{ x1, }\DataTypeTok{y =}\NormalTok{ y1)) }\OperatorTok{+}
\StringTok{    }\KeywordTok{geom\_point}\NormalTok{()}
\end{Highlighting}
\end{Shaded}

\includegraphics{book_files/figure-latex/unnamed-chunk-1-1.pdf}

\textbf{Exercise:} plot the other three datasets

\begin{Shaded}
\begin{Highlighting}[]
\NormalTok{anscombe}\OperatorTok{$}\NormalTok{x1}
\end{Highlighting}
\end{Shaded}

\begin{verbatim}
##  [1] 10  8 13  9 11 14  6  4 12  7  5
\end{verbatim}

\begin{Shaded}
\begin{Highlighting}[]
\NormalTok{anscombe[[}\StringTok{\textquotesingle{}x1\textquotesingle{}}\NormalTok{]]}
\end{Highlighting}
\end{Shaded}

\begin{verbatim}
##  [1] 10  8 13  9 11 14  6  4 12  7  5
\end{verbatim}

\begin{Shaded}
\begin{Highlighting}[]
\NormalTok{anscombe[,}\StringTok{\textquotesingle{}x1\textquotesingle{}}\NormalTok{]}
\end{Highlighting}
\end{Shaded}

\begin{verbatim}
##  [1] 10  8 13  9 11 14  6  4 12  7  5
\end{verbatim}

\begin{Shaded}
\begin{Highlighting}[]
\NormalTok{anscombe[,}\DecValTok{1}\NormalTok{]}
\end{Highlighting}
\end{Shaded}

\begin{verbatim}
##  [1] 10  8 13  9 11 14  6  4 12  7  5
\end{verbatim}

\begin{Shaded}
\begin{Highlighting}[]
\KeywordTok{mean}\NormalTok{(anscombe}\OperatorTok{$}\NormalTok{x1)}
\end{Highlighting}
\end{Shaded}

\begin{verbatim}
## [1] 9
\end{verbatim}

\begin{Shaded}
\begin{Highlighting}[]
\KeywordTok{mean}\NormalTok{(anscombe}\OperatorTok{$}\NormalTok{x2)}
\end{Highlighting}
\end{Shaded}

\begin{verbatim}
## [1] 9
\end{verbatim}

\begin{Shaded}
\begin{Highlighting}[]
\KeywordTok{sd}\NormalTok{(anscombe}\OperatorTok{$}\NormalTok{y1)}
\end{Highlighting}
\end{Shaded}

\begin{verbatim}
## [1] 2.031568
\end{verbatim}

\begin{Shaded}
\begin{Highlighting}[]
\KeywordTok{sd}\NormalTok{(anscombe}\OperatorTok{$}\NormalTok{y2)}
\end{Highlighting}
\end{Shaded}

\begin{verbatim}
## [1] 2.031657
\end{verbatim}

\begin{Shaded}
\begin{Highlighting}[]
\KeywordTok{cor}\NormalTok{(anscombe}\OperatorTok{$}\NormalTok{x1, anscombe}\OperatorTok{$}\NormalTok{y1)}
\end{Highlighting}
\end{Shaded}

\begin{verbatim}
## [1] 0.8164205
\end{verbatim}

\textbf{Exercise}: Calculate correlations for all four x-y pairs

\begin{itemize}
\tightlist
\item
  Some idiosyncracies of my code

  \begin{itemize}
  \tightlist
  \item
    \texttt{\textquotesingle{}} rather than \texttt{"}
  \item
    \texttt{\#\#} for documentation comment
  \item
    \texttt{=} not \texttt{\textless{}-} for assignment
  \end{itemize}
\item
  Essential Global Options

  \begin{itemize}
  \tightlist
  \item
    Never save .Rdata, load workspace
  \end{itemize}
\item
  Super useful custom keyboard commands:

  \begin{itemize}
  \tightlist
  \item
    Tools \(\to\) Modify Keyboard Shortcuts\ldots{}
  \item
    ⌘+1, ⌘+2, ⌘+3
  \end{itemize}
\end{itemize}

\hypertarget{base-r-data-types}{%
\section{Base R data types}\label{base-r-data-types}}

\hypertarget{numerics}{%
\subsection{Numerics}\label{numerics}}

\begin{Shaded}
\begin{Highlighting}[]
\NormalTok{foo =}\StringTok{ }\DecValTok{2}
\NormalTok{bar =}\StringTok{ }\DecValTok{3}
\NormalTok{foobar =}\StringTok{ }\NormalTok{foo }\OperatorTok{+}\StringTok{ }\NormalTok{bar}
\NormalTok{foo }\OperatorTok{/}\StringTok{ }\NormalTok{bar}
\end{Highlighting}
\end{Shaded}

\begin{verbatim}
## [1] 0.6666667
\end{verbatim}

\begin{Shaded}
\begin{Highlighting}[]
\KeywordTok{try}\NormalTok{(foo }\OperatorTok{/}\StringTok{ }\DecValTok{0}\NormalTok{)}
\end{Highlighting}
\end{Shaded}

\begin{verbatim}
## [1] Inf
\end{verbatim}

\hypertarget{vectors}{%
\subsection{Vectors}\label{vectors}}

\begin{Shaded}
\begin{Highlighting}[]
\NormalTok{foo =}\StringTok{ }\KeywordTok{c}\NormalTok{(}\DecValTok{2}\NormalTok{, }\DecValTok{3}\NormalTok{, }\DecValTok{4}\NormalTok{)}
\NormalTok{bar =}\StringTok{ }\KeywordTok{c}\NormalTok{(}\DecValTok{5}\NormalTok{, }\DecValTok{6}\NormalTok{, }\DecValTok{7}\NormalTok{)}
\NormalTok{foo[}\DecValTok{2}\NormalTok{]}
\end{Highlighting}
\end{Shaded}

\begin{verbatim}
## [1] 3
\end{verbatim}

\begin{Shaded}
\begin{Highlighting}[]
\NormalTok{foo[}\DecValTok{2}\OperatorTok{:}\DecValTok{3}\NormalTok{]}
\end{Highlighting}
\end{Shaded}

\begin{verbatim}
## [1] 3 4
\end{verbatim}

\begin{Shaded}
\begin{Highlighting}[]
\KeywordTok{try}\NormalTok{(foo[}\DecValTok{1}\NormalTok{, }\DecValTok{3}\NormalTok{])}
\end{Highlighting}
\end{Shaded}

\begin{verbatim}
## Error in foo[1, 3] : incorrect number of dimensions
\end{verbatim}

\begin{Shaded}
\begin{Highlighting}[]
\NormalTok{foo[}\KeywordTok{c}\NormalTok{(}\DecValTok{1}\NormalTok{, }\DecValTok{3}\NormalTok{)]}
\end{Highlighting}
\end{Shaded}

\begin{verbatim}
## [1] 2 4
\end{verbatim}

\begin{Shaded}
\begin{Highlighting}[]
\NormalTok{foo }\OperatorTok{+}\StringTok{ }\NormalTok{bar}
\end{Highlighting}
\end{Shaded}

\begin{verbatim}
## [1]  7  9 11
\end{verbatim}

\begin{Shaded}
\begin{Highlighting}[]
\NormalTok{foo }\OperatorTok{*}\StringTok{ }\NormalTok{bar}
\end{Highlighting}
\end{Shaded}

\begin{verbatim}
## [1] 10 18 28
\end{verbatim}

\begin{Shaded}
\begin{Highlighting}[]
\CommentTok{\#\# NB difference between vector and length{-}1 comparisons}
\NormalTok{bar }\OperatorTok{==}\StringTok{ }\NormalTok{foo }\OperatorTok{+}\StringTok{ }\DecValTok{3}
\end{Highlighting}
\end{Shaded}

\begin{verbatim}
## [1] TRUE TRUE TRUE
\end{verbatim}

\begin{Shaded}
\begin{Highlighting}[]
\KeywordTok{identical}\NormalTok{(bar, foo }\OperatorTok{+}\StringTok{ }\DecValTok{3}\NormalTok{)}
\end{Highlighting}
\end{Shaded}

\begin{verbatim}
## [1] TRUE
\end{verbatim}

\begin{Shaded}
\begin{Highlighting}[]
\KeywordTok{c}\NormalTok{(foo, bar)}
\end{Highlighting}
\end{Shaded}

\begin{verbatim}
## [1] 2 3 4 5 6 7
\end{verbatim}

\hypertarget{strings}{%
\subsection{Strings}\label{strings}}

\begin{Shaded}
\begin{Highlighting}[]
\NormalTok{foo =}\StringTok{ }\KeywordTok{c}\NormalTok{(}\StringTok{\textquotesingle{}monkey\textquotesingle{}}\NormalTok{, }\StringTok{\textquotesingle{}banana\textquotesingle{}}\NormalTok{)}
\NormalTok{bar =}\StringTok{ }\KeywordTok{c}\NormalTok{(}\StringTok{\textquotesingle{}house\textquotesingle{}}\NormalTok{)}
\CommentTok{\#\# How to put these together? }
\KeywordTok{try}\NormalTok{(foo }\OperatorTok{+}\StringTok{ }\NormalTok{bar)}
\end{Highlighting}
\end{Shaded}

\begin{verbatim}
## Error in foo + bar : non-numeric argument to binary operator
\end{verbatim}

\begin{Shaded}
\begin{Highlighting}[]
\KeywordTok{c}\NormalTok{(foo, bar)}
\end{Highlighting}
\end{Shaded}

\begin{verbatim}
## [1] "monkey" "banana" "house"
\end{verbatim}

\begin{Shaded}
\begin{Highlighting}[]
\KeywordTok{paste}\NormalTok{(foo, bar)}
\end{Highlighting}
\end{Shaded}

\begin{verbatim}
## [1] "monkey house" "banana house"
\end{verbatim}

\begin{Shaded}
\begin{Highlighting}[]
\KeywordTok{paste0}\NormalTok{(foo, bar)}
\end{Highlighting}
\end{Shaded}

\begin{verbatim}
## [1] "monkeyhouse" "bananahouse"
\end{verbatim}

\hypertarget{factors}{%
\subsection{Factors}\label{factors}}

\begin{Shaded}
\begin{Highlighting}[]
\NormalTok{foo\_str =}\StringTok{ }\KeywordTok{c}\NormalTok{(}\StringTok{\textquotesingle{}foo\textquotesingle{}}\NormalTok{, }\StringTok{\textquotesingle{}foo\textquotesingle{}}\NormalTok{, }\StringTok{\textquotesingle{}bar\textquotesingle{}}\NormalTok{, }\StringTok{\textquotesingle{}zoo\textquotesingle{}}\NormalTok{)}
\KeywordTok{as.integer}\NormalTok{(foo\_str)}
\end{Highlighting}
\end{Shaded}

\begin{verbatim}
## Warning: NAs introduced by coercion
\end{verbatim}

\begin{verbatim}
## [1] NA NA NA NA
\end{verbatim}

\begin{Shaded}
\begin{Highlighting}[]
\NormalTok{foo\_fct =}\StringTok{ }\KeywordTok{as.factor}\NormalTok{(foo\_str)}
\KeywordTok{as.integer}\NormalTok{(foo)}
\end{Highlighting}
\end{Shaded}

\begin{verbatim}
## Warning: NAs introduced by coercion
\end{verbatim}

\begin{verbatim}
## [1] NA NA
\end{verbatim}

\begin{Shaded}
\begin{Highlighting}[]
\KeywordTok{levels}\NormalTok{(foo\_fct)}
\end{Highlighting}
\end{Shaded}

\begin{verbatim}
## [1] "bar" "foo" "zoo"
\end{verbatim}

\hypertarget{lists}{%
\subsection{Lists}\label{lists}}

\emph{{[}todo{]}}

\hypertarget{data-frames}{%
\subsection{Data frames}\label{data-frames}}

\emph{{[}todo{]}}

\hypertarget{examining-variables}{%
\subsection{\texorpdfstring{\emph{{[}examining variables{]}}}{{[}examining variables{]}}}\label{examining-variables}}

\begin{itemize}
\tightlist
\item
  length
\item
  str
\item
  class, typeof
\item
  \href{https://gge-ucd.github.io/R-DAVIS/lesson_how_r_thinks_about_data.html\#challenge}{challenge}
\end{itemize}

\hypertarget{missing-data}{%
\subsection{\texorpdfstring{\emph{{[}Missing data{]}}}{{[}Missing data{]}}}\label{missing-data}}

\url{https://gge-ucd.github.io/R-DAVIS/lesson_how_r_thinks_about_data.html\#missing_data}

\hypertarget{git-and-data-management}{%
\chapter{Git and Data Management}\label{git-and-data-management}}

\hypertarget{cleaning-data}{%
\chapter{Cleaning Data}\label{cleaning-data}}

\hypertarget{errors-warnings-and-getting-help}{%
\chapter{Errors, Warnings, and Getting Help}\label{errors-warnings-and-getting-help}}

\hypertarget{exploratory-data-analysis}{%
\chapter{Exploratory Data Analysis}\label{exploratory-data-analysis}}

\printbibliography

\end{document}
